% Part: lambda-calculus
% Chapter: introduction
% Section: unique-readability

\documentclass[../../../include/open-logic-section]{subfiles}

\begin{document}

\olfileid[id]{lam}{syn}{unq}
\olsection{Keterbacaan Unik}

Kita mungkin bertanya-tanya apakah setiap suku hanya dapat dibentuk dengan
satu cara, dan memang demikian. Setiap suku lambda hanya dapat dibangun dan
ditafsirkan dengan satu cara. Dalam pembahasan berikut, suatu
\emph{pembentukan} adalah prosedur membangun sebuah suku dengan menggunakan
aturan-aturan pembentukan (satu kali atau beberapa kali) dari
\olref[trm]{defn:term}.

\begin{lem}\ollabel{lem:term-start}
Suatu suku diawali oleh variabel atau tanda kurung.
\end{lem}

\begin{proof}
Sesuatu tergolong suku hanya jika dibangun menurut
\olref[trm]{defn:term}. Jika sesuatu itu merupakan hasil
\olref[trm]{defn:term-var}, maka ia pasti sebuah variabel. Jika ia merupakan
hasil \olref[trm]{defn:term-abs} atau \olref[trm]{defn:term-app}, maka ia
diawali oleh tanda kurung.
\end{proof}

\begin{lem}\ollabel{lem:app-start}
Hasil suatu aplikasi diawali oleh dua tanda kurung atau oleh sebuah tanda
kurung dan sebuah variabel.
\end{lem}

\begin{proof}
Jika $M$ merupakan hasil suatu aplikasi, maka suku itu berbentuk $(PQ)$,
sehingga diawali oleh tanda kurung. Karena $P$ adalah suku, menurut
\olref{lem:term-start}, ia diawali oleh tanda kurung atau variabel.
\end{proof}

\begin{lem}\ollabel{lem:initial}
Tidak ada bagian awal sejati dari suatu suku yang juga merupakan suku.
\end{lem}

\begin{prob}
Buktikan \olref[lam][syn][unq]{lem:initial} dengan induksi pada panjang suku.
\end{prob}

\begin{prop}[Keterbacaan Unik] \ollabel{prop:unq}
Setiap suku mempunyai pembentukan yang unik. Dengan kata lain, jika suatu suku
$M$ dibentuk oleh sebuah pembentukan, maka itulah satu-satunya pembentukan yang
dapat membentuk suku tersebut.
\end{prop}

\begin{proof}
  Kita membuktikan hal ini dengan induksi pada pembentukan suku.
  \begin{enumerate}
    \item $M$ berbentuk
      $x$, dengan $x$ suatu variabel. Karena hasil abstraksi dan aplikasi
      selalu diawali oleh tanda kurung, keduanya tidak mungkin digunakan untuk
      membangun $M$. Jadi, pembentukan $M$ pasti merupakan satu langkah
      \olref[trm]{defn:term}\olref[trm]{defn:term-var}.
    \item $M$ berbentuk $(\lambd[x][N])$, dengan $x$ suatu variabel dan $N$
      suatu suku. Suku ini tidak mungkin dibangun menurut
      \olref[trm]{defn:term}\olref[trm]{defn:term-var}, karena bukan sebuah
      variabel tunggal. Menurut \olref{lem:app-start}, suku ini juga bukan
      hasil suatu aplikasi. Jadi, $M$ hanya mungkin merupakan hasil abstraksi
      pada~$N$. Berdasarkan hipotesis induksi, kita mengetahui bahwa
      pembentukan $N$ sendiri unik.
    \item $M$ berbentuk $(PQ)$, dengan $P$ dan $Q$ suku. Karena diawali oleh
      tanda kurung, suku ini tidak mungkin pula dibangun dengan
      \olref[trm]{defn:term}\olref[trm]{defn:term-var}. Menurut
      \olref{lem:term-start}, $P$ tidak mungkin diawali oleh $\lambd$, sehingga
      $(PQ)$ tidak mungkin merupakan hasil suatu abstraksi. Sekarang andaikan
      terdapat cara lain untuk membangun $M$ dengan aplikasi, misalnya suku
      itu juga berbentuk $(P'Q')$. Maka $P$ merupakan segmen awal sejati dari
      $P'$ (atau sebaliknya), dan hal ini mustahil menurut
      \olref{lem:initial}.
      Jadi, $P$ dan $Q$ ditentukan secara unik, dan berdasarkan hipotesis
      induksi kita mengetahui bahwa pembentukan $P$ dan $Q$ juga unik.
  \end{enumerate}
\end{proof}

Parafrasa yang lebih mudah dibaca dari proposisi di atas adalah sebagai
berikut:
\begin{prop}
  Suatu suku $M$ hanya dapat mempunyai salah satu bentuk berikut:
  \begin{enumerate}
    \item $x$, dengan $x$ suatu variabel yang ditentukan secara unik oleh $M$.
    \item $(\lambd[x][N])$, dengan $x$ suatu variabel dan $N$ suku lain;
      keduanya ditentukan secara unik oleh $M$.
    \item $(PQ)$, dengan $P$ dan $Q$ dua suku yang ditentukan secara unik
      oleh $M$.
  \end{enumerate}
\end{prop}

\end{document}
